% =============================================================================
% 012 / 068
% Status: raw
% =============================================================================
% Notes:
%
% Source question:
% 也不算問題吧。我只在網上學過國際音標和一點基礎語音學，而您的和其他人的吳語拼音教程沒有標註寬式國際音標，而我本地是最西部的太湖片地區與太湖片核心地區很多具體字音不一樣，所以遇到一些學習困難。另外，我覺得您應該推廣（或自行設計）更貼近英語拼讀規律的吳語拼音，比如“zh.wikipedia.org/wiki/上海話拉丁化方案”上的“江拉、滬一、錢一”。
% =============================================================================

\chapter[也不算問題吧。我只在網上學過國際音標和一點基礎語音學，而您的和其他人的吳語拼音教程沒有標註寬式國際音標，而我本地是最西部的太湖片地區與太湖片核心地區很多具體字音不一樣，所以遇到一些學習困難。另外，我覺得您應該推廣（或自行設計）更貼近英語拼讀規律的吳語拼音，比如“zh.wikipedia.org/wiki/上海話拉丁化方案”上的“江拉、滬一、錢一”。]{也不算問題吧。我只在網上學過國際音標和一點基礎語音學，而您的和其他人的吳語拼音教程沒有標註寬式國際音標，而我本地是最西部的太湖片地區與太湖片核心地區很多具體字音不一樣，所以遇到一些學習困難。另外，我覺得您應該推廣（或自行設計）更貼近英語拼讀規律的吳語拼音，比如“zh.wikipedia.org/wiki/上海話拉丁化方案”上的“江拉、滬一、錢一”。}
\qaanswer{
妳提出的是一個很多沒接觸過吳語拼音的朋友會問的典型問題

首先吳語拼音不是音值拼音，而是音位拼音，所謂音位指的是一個音位對應多個音值，比如同一個音位韻母eu，在吳語區各地的音值是不同的，這樣方便各地人按自己的習慣讀音，不需要人為的規定什麼標準音(定標準音又是一個誰也不服誰的口水仗)。

如果按國際音標來，那麼實際情況下吳語區一個鎮一個村就是一個拼法，整套系統下來估計會是天文數字吧！而且國際音標學習成本極高。事實上英語就是音位拼音語言，同一個拼寫英語區各地實際讀音都不太一樣。

吳語拼音是目前我個人覺得最科學也是最便於學習的一套拼音方案。希望我的回答能讓妳滿意。
}

